\movedhead{Behavioural biases in human investors}
\label{app:humaninvestors}
The biases this paper tests are not hypotheses about language models. They
are documented regularities in human investors, and the design borrows their
operationalisations rather than inventing new ones. Prospect theory
\citep{kahneman1979prospect} supplies the reference dependence that makes a
cost basis matter at all, and the framing manipulations of
\citet{tversky1981framing} supply the form of the framing family here: the
same information, valenced two ways. The disposition effect was named by
\citet{shefrin1985disposition} as a reluctance to realise losses and measured
in brokerage account data by \citet{odean1998losses}, who scores which
positions investors sell when they sell something. Mental accounting
\citep{thaler1985mental} motivates the sleeve manipulation, myopic loss
aversion \citep{benartzi1995myopic} the evaluation horizon, and home bias
\citep{french1991homebias} the domestic and foreign arms.
\citet{barber2000hazardous} supply the reason any of this is worth
measuring in an advisor: these biases carry a cost in realised returns.

These are also the source of the human baselines in
Section~\ref{sec:humanbaselines}, and the comparison is indicative rather
than like for like, because a field study of brokerage accounts and a
prompt battery are different instruments.
% Closed 2026-10-26. The Odean rates in Section 4.2 are read from Table I on
% p.1783 of the primary paper, not from any summary of it. Any further human
% rate added later must clear the same bar.
